\section{The CREST Benchmark}
\label{sec:design}

\paragraph{Construct.} We adopt the standard definition of creativity---an artifact must be both \emph{novel} and \emph{useful/appropriate} \citep{runco2012standard,sternberg1999concept}---augmented with Boden's third criterion, \emph{surprise} \citep{boden2004creative}, giving the New/Useful/Surprising (N/U/S) rubric---each axis scored $1$--$10$, with the \emph{graded} conjunction $\min(N,U,S)$ as the primary endpoint. A thresholded ``all three clear a fixed cutoff'' indicator was measured and \textbf{withdrawn} as primary: its rate swung $5.4$--$36.1\%$ across cutoffs $5$--$7$ on the pilot, so it is retained only for backward comparability. Crucially, \emph{Useful} functions as a factual-grounding check on \emph{New}, and \emph{Surprising} is expectation-violation \emph{conditioned on correctness}, which distinguishes CREST from novelty-only metrics (\S\ref{sec:related}). This N/U/S operationalization is the established construct---from the standard creativity literature and from recent deep-research-agent creativity evaluations \citep{caesar2026}---not a metric we introduce; it is defined independently of any evaluated agent.

\paragraph{Categories.} Creativity is not monolithic. CREST scores N/U/S across five categories that instantiate distinct creative processes from the literature: \textbf{C1 Constrained} and \textbf{C5 Open-Ended} correspond to Boden's \emph{exploratory} (search within a fixed space) and \emph{transformational} (rewriting the space) creativity; \textbf{C3 Cross-Domain} instantiates \emph{combinational} creativity, grounded in structure-mapping theory, in which cross-domain synthesis is relational alignment under the systematicity principle \citep{gentner1983structure,falkenhainer1989sme}; \textbf{C2 Counterfactual} is grounded in the rational imagination, where ``what-if'' reasoning is systematic over a causal model \citep{byrne2005rational}; \textbf{C4 Meta-Creativity} captures reflective, space-altering creativity.

\paragraph{Items.} CREST ships 50 expert-authored challenges---10 per category---across six domains (biology/medicine, materials/chemistry, economics/business, law/policy, engineering, arts/humanities; 8--9 each). Each challenge requires deep multi-source synthesis (never a single lookup) and ships with (a) a \emph{dated known-answers list} (the baseline against which \emph{New} is judged), (b) checkable \emph{hard rules} (the reference for \emph{Useful}), and (c) an \emph{expert baseline} answer (the reference for \emph{Surprising}). Every item is labeled by isolated axis, domain, and pre/post-training-cutoff.

\paragraph{Open prompts, judge-side rubric.} A design choice we state explicitly: the \emph{agent-facing} prompt is short and open (median $\sim$110 words) and poses only the creative question. All axis-isolation and anti-contamination machinery---the hard rules, dated known-answers, and expert baselines---lives \emph{judge-side} and is never shown to the agent. This preserves creative headroom and avoids over-steering the agent toward a target answer, while keeping grading reproducible and axis-isolated.

\paragraph{Contamination and length controls: what exists, and what does not.} Entrants run under their \emph{provider-native web search} (no separate search key), so all $250$ records come from \emph{live} retrieval and CREST is \textbf{not} contamination-controlled. The intended fair, reproducible run mode---a dated corpus snapshot queried through one shared retrieval index---is \textbf{not implemented}: \texttt{FrozenIndex} raises, so there is \textbf{no} frozen shared index and \textbf{no} released snapshot manifest, and no party can reconstruct the corpus any run actually saw. The one anti-contamination property that does hold is indirect: because \emph{New} is judged against per-item \emph{dated known-answers}, fresher retrieval cannot by itself inflate the axis. Four further controls named by earlier drafts do \textbf{not} exist. Canary strings for post-hoc training-set detection are \textbf{planned but absent} from all 50 items. Items carry a time-sensitivity label, but \textbf{no} released analysis splits results on training cutoff. There is \textbf{no} rolling refresh set; the release is a fixed 50 items. Output length is recorded per record but is \textbf{not} a covariate in any released analysis, so the well-documented length confound is unaddressed. And entrants do \textbf{not} receive a matched tool/compute budget: realised searches per item run from $0.64$ (\texttt{gemini\_shallow}, zero searches on $40$ of $50$ items) to $76.28$ (\texttt{gemini\_deepresearch\_max}, max $144$), so retrieval depth is confounded with vendor rather than controlled (\S\ref{sec:limitations}).

\paragraph{Release and reproducibility.} We release the questions, schema, judge configuration, rubric, analysis code and the promoted \texttt{result/leaderboard.json} under a documented license, with \textbf{Croissant} $+$ Croissant-RAI metadata and a \textbf{Datasheet for Datasets} \citep{gebru2021datasheets,croissant2024} covering provenance, an author-conflict-of-interest statement, and the adversarial item-vetting log. Two release artifacts named by earlier drafts are absent because the things they describe do not exist: there is \textbf{no} frozen-index manifest, and \textbf{no} human-study IRB or consent record, since no human study has been run. The release is a fixed 50 items, 10 per category, all 50 usable (no item was gated out of the reported run); there is \textbf{no} rolling refresh. The $250$-record run directory is \textbf{not} distributed, so the reported results cannot be recomputed from a fresh clone (\S\ref{sec:experiments}).
